% GENERATED FILE. Edit canonical lessons, then run npm run generate:books.
% canonical-source-hash: fnv1a64:f104655c7efbad8e
% canonical-lessons: JA-C181-kaimono, JA-C181-nedan, JA-C181-ikura, JA-C181-uru, JA-C181-ureru

\chapter{Prices and Selling}
\label{ch:ja-prices-and-selling}
\hlchaptermodality{181}

\begin{chapteropening}
\textbf{By the end of this chapter:} \emph{I can ask prices and say what a shop sells.}

Complete the last lesson of chapter 181: I can ask prices and say what a shop sells.
\end{chapteropening}

\section[\texorpdfstring{\ja{かいもの} (kaimono)}{かいもの (kaimono)}]{\ja{かいもの} (kaimono) — shopping; purchases}

\label{lesson:JA-C181-kaimono}



\begin{quote}
Before the new one: say the Japanese for to grip, then the Japanese for rice ball.
\end{quote}

\subsection*{You'll want to know: \ja{かいもの}}
\textbf{\ja{かいもの}} — \emph{kaimono} — \textquotedblleft{}shopping; purchases\textquotedblright{}.

\begin{quote}
\textbf{\ja{まちで} \ja{かいものを} \ja{します。}} — \emph{machi de kaimono o shimasu} — \textquotedblleft{}I go shopping in town.\textquotedblright{}
\end{quote}

\begin{cousinweb}
\textbf{\ja{かいもの}} (\emph{kaimono}) is native Japanese, built from \emph{kai}, the stem of \emph{kau}, \textquotedblleft{}to buy\textquotedblright{}, and \emph{mono}, \textquotedblleft{}thing\textquotedblright{} (JMdict, KANJIDIC2).
\end{cousinweb}

\begin{grammarlens}[title={What you've built}]
One more word for everyday life at A2.
\end{grammarlens}

\subsection*{Your turn}
\begin{itemize}
\raggedright
  \item \emph{Say it:} \emph{kaimono}
  \item \emph{Say it:} \emph{kaimono}, once more
  \item \emph{Say it:} \emph{machi de kaimono o shimasu} — I go shopping in town.
  \item \emph{Recall:} say the Japanese for to grip, then the Japanese for rice ball, then say \emph{kaimono} again
\end{itemize}

\begin{tcolorbox}[breakable,colback=teal!4,colframe=teal!35!black,title={Before you move on}]
What does \ja{かいもの} mean? (\textquotedblleft{}Shopping; purchases\textquotedblright{}.) Say it once more.
\end{tcolorbox}
\section[\texorpdfstring{\ja{ねだん} (nedan)}{ねだん (nedan)}]{\ja{ねだん} (nedan) — price; cost}

\label{lesson:JA-C181-nedan}



\begin{quote}
Before the new one: say the Japanese for rice ball, then the Japanese for shopping.
\end{quote}

\subsection*{You'll want to know: \ja{ねだん}}
\textbf{\ja{ねだん}} — \emph{nedan} — \textquotedblleft{}price; cost\textquotedblright{}.

\begin{quote}
\textbf{\ja{ねだんは} \ja{せんえんです。}} — \emph{nedan wa sen'en desu} — \textquotedblleft{}The price is a thousand yen.\textquotedblright{}
\end{quote}

\begin{cousinweb}
\textbf{\ja{ねだん}} (\emph{nedan}) is built from native \emph{ne}, \textquotedblleft{}price, value\textquotedblright{}, and Sino-Japanese \emph{dan}, \textquotedblleft{}grade, step\textquotedblright{}, a reading borrowed from Chinese (Middle Chinese \emph{dwanH}) (JMdict, KANJIDIC2).
\end{cousinweb}

\begin{grammarlens}[title={What you've built}]
One more word for everyday life at A2.
\end{grammarlens}

\subsection*{Your turn}
\begin{itemize}
\raggedright
  \item \emph{Say it:} \emph{nedan}
  \item \emph{Say it:} \emph{nedan}, once more
  \item \emph{Say it:} \emph{nedan wa sen'en desu} — The price is a thousand yen.
  \item \emph{Recall:} say the Japanese for rice ball, then the Japanese for shopping, then say \emph{nedan} again
\end{itemize}

\begin{tcolorbox}[breakable,colback=teal!4,colframe=teal!35!black,title={Before you move on}]
What does \ja{ねだん} mean? (\textquotedblleft{}Price; cost\textquotedblright{}.) Say it once more.
\end{tcolorbox}
\section[\texorpdfstring{\ja{いくら} (ikura)}{いくら (ikura)}]{\ja{いくら} (ikura) — how much}

\label{lesson:JA-C181-ikura}



\begin{quote}
Before the new one: say the Japanese for shopping, then the Japanese for price.
\end{quote}

\subsection*{You'll want to know: \ja{いくら}}
\textbf{\ja{いくら}} — \emph{ikura} — \textquotedblleft{}how much\textquotedblright{}.

\begin{quote}
\textbf{\ja{この} \ja{ほんは} \ja{いくらですか。}} — \emph{kono hon wa ikura desu ka} — \textquotedblleft{}How much is this book?\textquotedblright{}
\end{quote}

\begin{cousinweb}
\textbf{\ja{いくら}} (\emph{ikura}) is a native Japanese word (wago), inherited rather than borrowed from Chinese; its kanji was attached later as a kun'yomi; JMdict marks it as usually written in kana (JMdict, KANJIDIC2).
\end{cousinweb}

\begin{grammarlens}[title={What you've built}]
One more word for everyday life at A2.
\end{grammarlens}

\subsection*{Your turn}
\begin{itemize}
\raggedright
  \item \emph{Say it:} \emph{ikura}
  \item \emph{Say it:} \emph{ikura}, once more
  \item \emph{Say it:} \emph{kono hon wa ikura desu ka} — How much is this book?
  \item \emph{Recall:} say the Japanese for shopping, then the Japanese for price, then say \emph{ikura} again
\end{itemize}

\begin{tcolorbox}[breakable,colback=teal!4,colframe=teal!35!black,title={Before you move on}]
What does \ja{いくら} mean? (\textquotedblleft{}How much\textquotedblright{}.) Say it once more.
\end{tcolorbox}
\section[\texorpdfstring{\ja{うる} (uru)}{うる (uru)}]{\ja{うる} (uru) — to sell}

\label{lesson:JA-C181-uru}



\begin{quote}
Before the new one: say the Japanese for price, then the Japanese for how much.
\end{quote}

\subsection*{You'll want to know: \ja{うる}}
\textbf{\ja{うる}} — \emph{uru} — \textquotedblleft{}to sell\textquotedblright{}.

\begin{quote}
\textbf{\ja{あの} \ja{みせは} \ja{やさいを} \ja{うります。}} — \emph{ano mise wa yasai o urimasu} — \textquotedblleft{}That shop sells vegetables.\textquotedblright{}
\end{quote}

Polite \textbf{\ja{うります}}, \emph{urimasu}; a godan verb.

\begin{cousinweb}
\textbf{\ja{うる}} (\emph{uru}) is a native Japanese word (wago), inherited rather than borrowed from Chinese; its kanji was attached later as a kun'yomi (JMdict, KANJIDIC2).
\end{cousinweb}

\begin{grammarlens}[title={What you've built}]
One more word for everyday life at A2.
\end{grammarlens}

\subsection*{Your turn}
\begin{itemize}
\raggedright
  \item \emph{Say it:} \emph{uru}
  \item \emph{Say it:} \emph{uru}, once more
  \item \emph{Say it:} \emph{ano mise wa yasai o urimasu} — That shop sells vegetables.
  \item \emph{Recall:} say the Japanese for price, then the Japanese for how much, then say \emph{uru} again
\end{itemize}

\begin{tcolorbox}[breakable,colback=teal!4,colframe=teal!35!black,title={Before you move on}]
What does \ja{うる} mean? (\textquotedblleft{}To sell\textquotedblright{}.) Say it once more.
\end{tcolorbox}
\section[\texorpdfstring{\ja{うれる} (ureru)}{うれる (ureru)}]{\ja{うれる} (ureru) — to sell well; to be popular}

\label{lesson:JA-C181-ureru}



\begin{quote}
Before the new one: say the Japanese for how much, then the Japanese for to sell.
\end{quote}

\subsection*{You'll want to know: \ja{うれる}}
\textbf{\ja{うれる}} — \emph{ureru} — \textquotedblleft{}to sell well; to be popular\textquotedblright{}.

\begin{quote}
\textbf{\ja{この} \ja{ほんは} \ja{よく} \ja{うれます。}} — \emph{kono hon wa yoku uremasu} — \textquotedblleft{}This book sells well.\textquotedblright{}
\end{quote}

Polite \textbf{\ja{うれます}}, \emph{uremasu}; an ichidan verb. Its transitive partner is \textbf{\ja{うる}}, \textquotedblleft{}to sell\textquotedblright{}.

\begin{cousinweb}
\textbf{\ja{うれる}} (\emph{ureru}) is a native Japanese word, inherited rather than borrowed, built on the same stem as \emph{uru}, \textquotedblleft{}to sell\textquotedblright{}; its kanji was attached later (JMdict, KANJIDIC2).
\end{cousinweb}

\begin{grammarlens}[title={What you've built}]
One more word for everyday life at A2.
\end{grammarlens}

\subsection*{Your turn}
\begin{itemize}
\raggedright
  \item \emph{Say it:} \emph{ureru}
  \item \emph{Say it:} \emph{ureru}, once more
  \item \emph{Say it:} \emph{kono hon wa yoku uremasu} — This book sells well.
  \item \emph{Recall:} say the Japanese for how much, then the Japanese for to sell, then say \emph{ureru} again
\end{itemize}

\begin{tcolorbox}[breakable,colback=teal!4,colframe=teal!35!black,title={Before you move on}]
What does \ja{うれる} mean? (\textquotedblleft{}To sell well; to be popular\textquotedblright{}.) Say it once more.
\end{tcolorbox}
